\section*{第 \ref{ch:b3:representations} 章——有限群的
表示}

\begin{solution}{exo:b3:representations:1}
(a) $\Z/n\Z$ 是交换群：所有\omterm{def:b3:representations:rep}{不可约}表示的次数均为 $1$
（\cref{prop:b3:representations:onedim}），即为同态
$\chi \colon \Z/n\Z \to \C^\times$，由 $\chi(\bar 1)
= \omega$ 唯一确定，且 $\omega^n = 1$：这 $n$ 个\omterm{def:b3:representations:character}{特征标} $\chi_k(\bar
m) = \eu^{2\iu\pi km/n}$。当 $n = 4$ 时（\omterm{ex:b3:groups:actions}{共轭类} $=$ 元素
$\bar0,\bar1,\bar2,\bar3$):
\[
\begin{array}{c|cccc}
& \bar0 & \bar1 & \bar2 & \bar3\\
\hline
\chi_0 & 1 & 1 & 1 & 1\\
\chi_1 & 1 & \iu & -1 & -\iu\\
\chi_2 & 1 & -1 & 1 & -1\\
\chi_3 & 1 & -\iu & -1 & \iu
\end{array}
\]
(b) 行：$\langle\chi_k,\chi_l\rangle = \frac14\sum_m
\eu^{2\iu\pi(l-k)m/4} = \delta_{kl}$（等比数列求和）；列的情形
同理。矩阵 $(\eu^{2\iu\pi km/n})_{k,m}$ 即
\emph{离散傅里叶变换}矩阵——与
本科二年级分册生成函数一章中所用的单位根滤子相同；
\omterm{def:b3:representations:character}{特征标}的正交性推广了
离散傅里叶变换的反演公式。
\end{solution}

\begin{solution}{exo:b3:representations:2}
(a) $\rho(g)$ 在基 $(e_x)$ 下的矩阵是一个
置换矩阵，其迹为 $x$ 中满足 $g\cdot x =
x$ 者的个数。于是
\[
\langle\mathbf 1, \chi\rangle = \frac1{\abs G}\sum_g
\abs{\operatorname{Fix}(g)} = \#\{\text{轨道数}\}
\]
由伯恩赛德计数引理（\cref{exo:b3:groups:5}）可得——
等价地，这也就是计算平凡
表示的重数，其同型子空间即
$G$-不变向量构成的空间，其维数等于\omterm{def:b3:groups:action}{轨道}的个数（每个
\omterm{def:b3:groups:action}{轨道}对应一个指示函数）。

(b) 由于 $\operatorname{Fix}_{X\times X}(g) =
\operatorname{Fix}_X(g)^2$，对 $X \times X$ 应用 (a)，
得 $\langle\chi,\chi\rangle = \frac1{\abs
G}\sum\abs{\operatorname{Fix}(g)}^2 = \#\{\text{ 在 }
X\times X\}$ 上的\omterm{def:b3:groups:action}{轨道}（$\chi$ 为实值）。记 $\chi = \mathbf 1 +
\psi$：$\langle\mathbf1,\chi\rangle = 1$（由于传递性），故
$\langle\psi,\psi\rangle = \langle\chi,\chi\rangle - 1$。
$X\times X$ 上的作用，以对角线为一个\omterm{def:b3:groups:action}{轨道}；
恰有一个另外的\omterm{def:b3:groups:action}{轨道}，当且仅当 $G$ 在相异数对上可迁：
$\langle\psi,\psi\rangle = 1$ 当且仅当 $2$-传递
（\cref{cor:b3:representations:multiplicity}）。

(c) $S_n$ 是 $2$-传递的（作用于 $\{1,\dots,n\}$ 上，任一
相异数对都可以映到任一相异数对）：$\psi = \chi_{\mathrm{std}}$ 是
\omterm{def:b3:representations:rep}{不可约}的。
\end{solution}

\begin{solution}{exo:b3:representations:3}
$S_3$：三个\omterm{ex:b3:groups:actions}{共轭类}，$\sum n_i^2 = 6 = 1 + 1 + 4$；两个
次数为 $1$ 的\omterm{def:b3:representations:character}{特征标}是 $\mathbf 1, \varepsilon$（$[S_3 :
A_3] = 2$）；第三行 $(2, a, b)$ 由
与 $e$ 列的列正交性：$1 - 1 + 2a = 0$ 且 $1 + 1 +
2b = 0$：$a = 0$，$b = -1$——即
\cref{ex:b3:representations:s3}中的表。

$S_4$ 的检验（行）：$\langle\chi_{\mathrm{std}},
\varepsilon\chi_{\mathrm{std}}\rangle = \frac1{24}(9 - 6 + 3 + 0
- 6) = 0$; $\langle\chi_2,\chi_2\rangle = \frac1{24}(4 + 0 + 12
+ 8 + 0) = 1$。列的检验：$e$ 与 $(1\,2)$：$1 - 1 + 0 + 3 - 3
= 0$；$(1\,2)$ 与自身：$1 + 1 + 0 + 1 + 1 = 4 =
\abs{Z_{S_4}((1\,2))} = 24/6$.

$4$ 点上的置换\omterm{def:b3:representations:character}{特征标}：$(4, 2, 0, 1, 0) = \mathbf
1 + \chi_{\mathrm{std}}$（不动点个数；减去首
行）。对 $\chi_{\mathrm{std}}^2 = (9, 1, 1, 0, 1)$：
\[
\langle\mathbf1,\cdot\rangle = \tfrac{9 + 6 + 3 + 0 + 6}{24} =
1,\quad
\langle\varepsilon,\cdot\rangle = 0,\quad
\langle\chi_2,\cdot\rangle = \tfrac{18 + 6}{24} = 1,\quad
\langle\chi_{\mathrm{std}},\cdot\rangle = 1,\quad
\langle\varepsilon\chi_{\mathrm{std}},\cdot\rangle = 1:
\]
$\chi_{\mathrm{std}}^2 = \mathbf 1 + \chi_2 +
\chi_{\mathrm{std}} + \varepsilon\chi_{\mathrm{std}}$
（维数：$9 = 1 + 2 + 3 + 3$）。
\end{solution}

\begin{solution}{exo:b3:representations:4}
两个群都有五个\omterm{ex:b3:groups:actions}{共轭类}，次数模式为 $(1,1,1,1,2)$
（四个次数为 $1$ 的\omterm{def:b3:representations:character}{特征标}，来自阿贝尔化 $\cong (\Z/2\Z)^2$，
再由 $\sum n_i^2 = 8$）。\omterm{ex:b3:groups:actions}{共轭类}依次取 $e$、$z$（
中心对合，分别为 $r^2$、$-1$），以及三个
两元素的\omterm{ex:b3:groups:actions}{共轭类}：
\[
\begin{array}{c|ccccc}
& e & z & C_1 & C_2 & C_3\\
\hline
\chi^{(1)} & 1 & 1 & 1 & 1 & 1\\
\chi^{(2)} & 1 & 1 & 1 & -1 & -1\\
\chi^{(3)} & 1 & 1 & -1 & 1 & -1\\
\chi^{(4)} & 1 & 1 & -1 & -1 & 1\\
\chi^{(5)} & 2 & -2 & 0 & 0 & 0
\end{array}
\]
（最后一行由列正交性得出）。
$D_4$ 与 $Q_8$ 的表完全相同，二者却不同构
（\cref{pb:b3:groups:1}）。该表\emph{确实}能刻画：$\abs
G$、\omterm{ex:b3:groups:actions}{共轭类}大小、中心（$\{g : \abs{\chi_i(g)} = n_i\
\forall i\}$：两者的阶均为 $2$）、阿贝尔化、整个
\omterm{def:b3:groups:normal}{正规子群}格（核与交，
\cref{exo:b3:representations:6}）。它无法刻画元素
阶：$D_4$ 有五个对合，$Q_8$ 只有一个——因此
同构类型确实比\omterm{def:b3:representations:table}{特征标表}更精细。
\end{solution}

\begin{solution}{exo:b3:representations:5}
(a) 在 $S_4$ 中，$(1\,2\,3)$ 的\omterm{ex:b3:groups:actions}{中心化子}阶为 $24/8 =
3$，即 $\langle(1\,2\,3)\rangle \subseteq A_4$。于是
$Z_{A_4}((1\,2\,3))$ 的阶为 $3$，而 $A_4$-共轭类含 $12/3
= 4$ 个元素：八个 $3$-轮换分裂为两个 $A_4$-共轭类
（代表元分别为 $(1\,2\,3)$ 及其逆）。

(b) $A_4/V \cong \Z/3\Z$ 给出三个次数为 $1$ 的\omterm{def:b3:representations:character}{特征标}
（$\omega = \eu^{2\iu\pi/3}$；$3$-轮换所在的类映为
$\bar1, \bar2$）；$\chi_{\mathrm{std}}$ 的限制
仍\omterm{def:b3:representations:rep}{不可约}（$\langle\chi,\chi\rangle = \frac1{12}(9 + 3
\cdot 1 + 0 + 0) = 1$）：
\[
\begin{array}{c|cccc}
A_4 & e\,[1] & (1\,2)(3\,4)\,[3] & (1\,2\,3)\,[4] &
(1\,3\,2)\,[4]\\
\hline
\mathbf 1 & 1 & 1 & 1 & 1\\
\chi_\omega & 1 & 1 & \omega & \omega^2\\
\chi_{\bar\omega} & 1 & 1 & \omega^2 & \omega\\
\chi_3 & 3 & -1 & 0 & 0
\end{array}
\]
(c) 各核：$\ker\chi_\omega = \ker\chi_{\bar\omega} = V$；
$\ker\chi_3 = \{e\}$（表中其余各项的模均不为 $3$）。\omterm{def:b3:groups:normal}{正规子群}
即为诸核之交
（\cref{exo:b3:representations:6}）：$\{e\}$、$V$、$A_4$——特别地，
$A_4$ 既无阶为 $2$ 的\omterm{def:b3:groups:normal}{正规子群}，
也无指数为 $2$ 的\omterm{def:b3:groups:normal}{正规子群}。
\end{solution}

\begin{solution}{exo:b3:representations:6}
(a) 若 $\chi(g) = \chi(e) = n$，则在
$\abs{\chi(g)} \leq n$ 中取等号，迫使 $\rho(g) = \lambda\,\mathrm{id}$
（\cref{prop:b3:representations:charbasics}），且 $n\lambda = n$：
$\rho(g) = \mathrm{id}$。逆命题显然成立。

(b) 设 $N \trianglelefteq G$，\omterm{thm:b3:groups:quotient}{商群}的\omterm{ex:b3:representations:examples}{正则表示}
$G/N$ 是忠实的；将其分解为 $G/N$ 的\omterm{def:b3:representations:rep}{不可约}表示，并
将它们提升到 $G$ 上（\cref{prop:b3:representations:onedim}(b)）：
$\chi_{i_1}, \dots$ 为 $G$ 的\omterm{def:b3:representations:rep}{不可约}\omterm{def:b3:representations:character}{特征标}，其核
包含 $N$，而其\emph{公共}核恰为
$\{e\}$ 的原像，即 $N$（由\omterm{thm:b3:groups:quotient}{商群}上的忠实性）。
于是 $N = \bigcap_j \ker\chi_{i_j}$。

(c) 若 $G$ 为\omterm{def:b3:groups:simple}{单群}：对非平凡的\omterm{def:b3:representations:rep}{不可约} $\chi$，
$\ker\chi \trianglelefteq G$ 不等于 $G$（\omterm{def:b3:representations:rep}{不可约}
表示若在整个 $G$ 上平凡，则即为平凡\omterm{def:b3:representations:character}{特征标}），
故 $\ker\chi = \{e\}$。反之，设所有非平凡
的核均平凡，且 $N \trianglelefteq G$ 满足 $N \neq
G$。在（b）中把 $N$ 表为核之交时，
其中某个所涉\omterm{def:b3:representations:character}{特征标}非平凡（若全部平凡，则
交集即为 $G$），其核为 $\{e\}$：$N =
\{e\}$。因此\omterm{def:b3:groups:normal}{正规子群}仅为 $\{e\}$ 与 $G$。
\end{solution}

\begin{solution}{exo:b3:representations:7}
$8$ 阶非阿贝尔群：次数为 $1$ 的\omterm{def:b3:representations:character}{特征标}个数为
$[G : D(G)]$，它是 $8$ 的真因数（非阿贝尔时 $D(G) \neq
\{e\}$），且 $\sum n_i^2 = 8$。记 $k$ 个次数为一，其余次数
$\geq 2$：$8 - k \equiv 0$，且各平方 $\geq 4$，又 $k
\mid 8$，$k < 8$。$k = 4$：恰一个次数为 $2$，相容。$k =
2$：剩下 $6$，它不是若干个平方 $\geq 4$ 之和。$k = 1$：
不可能，因 $k = [G : D(G)] \geq 2$——$G/D(G)$ 为
非平凡的阿贝尔 $2$-群，因 $G$ 是 $2$-群且
$D(G) \neq G$，由 $p$-群的可解性
（\cref{ex:b3:groups:solvableexamples}）。故次数组为
$(1,1,1,1,2)$。$\chi_5$ 的忠实性：四个次数为 $1$ 的
\omterm{def:b3:representations:character}{特征标}之核均含 $D(G)$；若 $\ker\chi_5
\supseteq N \neq \{e\}$ 对某个极小\omterm{def:b3:groups:normal}{正规子群} $N$（$N
\subseteq D(G)$ 与否皆可，取 $N \subseteq \ker\chi_5$
非平凡），则 $N$ 将落在全部五个核中，这些核的
交为平凡（\omterm{ex:b3:representations:examples}{正则表示}
是忠实的）：矛盾。$p^3$ 阶非阿贝尔群中，$Z(G)$ 的阶
为 $p$（若阶为 $p^2$，则 $G/Z$ 为循环群，$G$ 为阿贝尔群），
$G/Z(G)$ 为阶 $p^2$ 的阿贝尔群，故 $D(G) \subseteq Z(G)$，
且 $D(G) \ne \{e\}$：$D(G) = Z(G)$，由此得 $p^2$ 个
次数为 $1$ 的\omterm{def:b3:representations:character}{特征标}。其余次数满足 $\sum n_i^2 = p^3 -
p^2$，且每个 $n_i > 1$ 整除 $\abs G$
（\cref{pb:b3:representations:1}，第 8 题），故 $n_i \in
\{p\}$（若 $n_i = p^2$ 则超出：$p^4 > p^3 - p^2$）：恰有 $p
- 1$ 个次数为 $p$ 的\omterm{def:b3:representations:character}{特征标}。
\end{solution}

\begin{solution}{exo:b3:representations:8}
对表示 $\rho, \sigma$，即 $G, H$ 分别在 $V, W$ 上的表示，定义
$\rho \boxtimes \sigma$ 为 $G \times H$ 的表示，作用于
$V \otimes W$ 上——具体地，在矩阵层面：$(\rho\boxtimes
\sigma)(g,h)$ 是克罗内克积 $\rho(g)\otimes
\sigma(h)$，其迹为
$\operatorname{tr}\rho(g)\operatorname{tr}\sigma(h) =
\chi(g)\psi(h)$（矩阵的克罗内克积 $A \otimes B$
的迹为 $\operatorname{tr}A\operatorname{tr}B$：其对角线
元素为 $a_{kk}b_{ll}$）。故 $\chi\psi$ 是\omterm{def:b3:representations:character}{特征标}，且
\[
\langle\chi\psi, \chi'\psi'\rangle_{G\times H}
= \frac{1}{\abs G\abs H}\sum_{g,h}
\overline{\chi(g)\psi(h)}\,\chi'(g)\psi'(h)
= \langle\chi,\chi'\rangle_G\,\langle\psi,\psi'\rangle_H
= \delta_{\chi\chi'}\delta_{\psi\psi'}.
\]
特别地，$\langle\chi\psi,\chi\psi\rangle = 1$，故每个
$\chi\psi$ \omterm{def:b3:representations:rep}{不可约}
（\cref{cor:b3:representations:multiplicity}）。它们是
$r_Gr_H$ 个互异的\omterm{def:b3:representations:rep}{不可约}\omterm{def:b3:representations:character}{特征标}；
$G\times H$ 的\omterm{ex:b3:groups:actions}{共轭类}即各类之积（$(g,h) \sim (g',h')$
按分量），故共 $r_Gr_H$ 个：列表
已完整（\cref{thm:b3:representations:numberirr}）。对
$(\Z/2\Z)^2$：四个符号\omterm{def:b3:representations:character}{特征标}
$(\pm1)\otimes(\pm1)$——正是
\cref{exo:b3:representations:4} 左上角块中的表。有限
阿贝尔群为循环群之积
（\cref{cor:b3:modules:abelian}）；其\omterm{def:b3:representations:rep}{不可约}\omterm{def:b3:representations:character}{特征标}为
各循环因子\omterm{def:b3:representations:character}{特征标}的乘积（\cref{exo:b3:representations:1}）：
次数均为 $1$。
\end{solution}

\begin{solution}{exo:b3:representations:9}
(a) $D(A_5) = A_5$（非阿贝尔\omterm{def:b3:groups:simple}{单群}），故次数为 $1$ 的
唯一\omterm{def:b3:representations:character}{特征标}为 $\mathbf 1$
（见\cref{prop:b3:representations:onedim}）。需满足 $n_2^2 + n_3^2
+ n_4^2 + n_5^2 = 59$，且每个 $n_i \geq 2$；逐一试验
平方数 $4, 9, 16, 25, 36, 49$，唯一可行的多重集为
$\{9, 9, 16, 25\}$：最大平方数为 $49$ 时，余数 $10$
不能写成三个平方数之和（每个 $\geq 4$）；取 $36$ 时，余数
$23$ 同样不行（$16 + 4 + 4 = 24$，$9 + 9 + 4 = 22$）；取
最大为 $25$ 时，验证 $25 + 16 + 9 + 9 = 59$ 成立，而 $25 +
25$、$25 + 16 + 16$、$25 + 16 + 4$ 等变式均不成立；取最大为
$16$ 时：$16\cdot3 = 48 < 59 - 4$。次数依次为 $1, 3, 3, 4, 5$。

(b) $5$ 个点上的置换，不动点数为 $(5, 1, 2, 0, 0)$，
于是 $\chi_4 = (4, 0, 1, -1, -1)$，且
$\langle\chi_4,\chi_4\rangle = \frac{16 + 0 + 20 + 12 + 12}{60}
= 1$，故\omterm{def:b3:representations:rep}{不可约}。作用于六个西罗 $5$-子群上，一个
对合元恰固定 $2$ 个（其\omterm{ex:b3:groups:actions}{正规化子}为\omterm{ex:b3:groups:semidirectexamples}{二面体群}，
阶为 $10$，各含 $5$ 个对合元：$30$ 个关联
对应 $15$ 个对合元），一个 $3$ 阶元固定 $0$ 个（$3$ 阶元在
$D_5$ 中不存在），一个 $5$ 阶元恰固定 $1$ 个（它恰含于唯一的
\omterm{def:b3:groups:sylow}{西罗子群}中）：置换\omterm{def:b3:representations:character}{特征标}为 $(6, 2, 0, 1, 1)$，且 $\chi_5 =
(5, 1, -1, 0, 0)$，其范数 $\frac{25 + 15 + 20 + 0 + 0}{60} =
1$，故\omterm{def:b3:representations:rep}{不可约}。余下两行 $(3, a, b, c, d)$、$(3, a', b', c',
d')$ 尚待确定。列范数
（见\cref{cor:b3:representations:column}）。在
$(1\,2)(3\,4)$ 所在的类上，$\abs{Z} = 4$：$1 + 0 + 1 + a^2 + a'^2 = 4$；
与 $e$ 列的正交性：$1\cdot1 + 4\cdot 0 + 5\cdot1 + 3(a + a') =
0$：$a + a' = -2$，$a^2 + a'^2 = 2$，故 $a = a' = -1$。在
$3$-轮换类上，$\abs Z = 3$：$1 + 1 + 1 + b^2 + b'^2 = 3$：$b = b'
= 0$。在每个 $5$-类上，$\abs Z = 5$：与 $e$ 列的正交性
给出 $1\cdot 1 + 4\cdot(-1) + 5\cdot 0 + 3(c + c') = 0$，故
$c + c' = 1$；且
$c^2 + c'^2 = 3$：$\{c, c'\} = \bigl\{\frac{1+\sqrt5}2,
\frac{1-\sqrt5}2\bigr\}$——黄金比及其
共轭数；第二个 $5$-类取互换后的值（
两行须正交）。

(c) 在完整表中，非平凡行的任何元素均不等于
其次数（第一列除外）：每个核 $\{g :
\chi_i(g) = n_i\}$ 均为平凡群。由
\cref{exo:b3:representations:6}(c) 可知 $A_5$ 为\omterm{def:b3:groups:simple}{单群}。
\end{solution}

\begin{solution}{exo:b3:representations:10}
$\rho(z)$ 与每个 $\rho(g)$ 可交换（$z$ 为中心元），即\
$\rho(z) \in \operatorname{End}_G(V) = \C\,\mathrm{id}$（舒尔）：
$\rho(z) = \lambda_z\,\mathrm{id}$，且 $z \mapsto \lambda_z$ 是
乘法的：$Z(G) \to \C^\times$ 为同态。则
$\chi(z) = n\lambda_z$，其中 $\abs{\lambda_z} = 1$（单位根
）：$\abs{\chi(z)} = n$。若 $\rho$ 为忠实表示，则 $\lambda$ 在
$Z(G)$ 上为单射（$\rho(z) = \mathrm{id} \iff \lambda_z =
1$），故 $Z(G)$ 嵌入 $\C^\times$；域的
乘法群的有限子群必为循环群
（见\cref{thm:b3:galois:cyclic}）。
\end{solution}

\begin{solution}{exo:b3:representations:11}
(a) 等变性：$\rho(h)p_i\rho(h)^{-1}$ 对求和作重新编号
（$g \mapsto hgh^{-1}$，且 $\chi_i$ 为\omterm{def:b3:representations:character}{类函数}）：
$p_i$ 与\omterm{def:b3:groups:action}{群作用}可交换。对\omterm{def:b3:representations:rep}{不可约}表示 $W$，
其\omterm{def:b3:representations:character}{特征标}为 $\chi_j$，由舒尔引理，其限制为一个数量
$\lambda\,\mathrm{id}_W$；取迹得
\[
\lambda\,n_j = \frac{n_i}{\abs G}\sum_g
\overline{\chi_i(g)}\,\chi_j(g) = n_i\,\langle\chi_i,
\chi_j\rangle = n_i\,\delta_{ij}
\]
（第一正交关系）：$\lambda = \delta_{ij}$。

(b) 将 $V$ 分解为\omterm{def:b3:representations:rep}{不可约}表示的直和（马施克）：$p_i$ 作用为
在\omterm{def:b3:representations:character}{特征标}为 $\chi_i$ 的直和项上为恒等，在
其余所有直和项上为 $0$，故 $p_i$ 是到它们之和 $V_i$
沿其余直和项之和；像 $V_i$ 不依赖于
所取的分解（它是被该投影所固定的向量全体，
$p_i$ 所定义，无需选取）。$\sum_ip_i$ 在
每个\omterm{def:b3:representations:rep}{不可约}直和项上为恒等，即 $\mathrm{id}_V$。
$V_i \cong W_i^{\oplus m_i}$ 的更细分解涉及
选取 $\operatorname{Hom}_G(W_i, V)$ 的一组基：典范的
则不然。

(c) 对 $\varepsilon$（次数为 $1$）：$p_\varepsilon =
\frac1{6}\sum_{g}\varepsilon(g)\,\rho(g)$，即群
代数中
\[
p_\varepsilon = \tfrac16\bigl(e - (1\,2) - (1\,3) - (2\,3)
+ (1\,2\,3) + (1\,3\,2)\bigr) .
\]
平方：$g$ 在 $p_\varepsilon^2$ 中的系数为
$\frac1{36}\sum_{h}\varepsilon(h)\varepsilon(h^{-1}g) =
\frac1{36}\,\varepsilon(g)\sum_h\varepsilon(h)^2 =
\frac{6}{36}\varepsilon(g)$: $p_\varepsilon^2 =
p_\varepsilon$。（其在\omterm{ex:b3:representations:examples}{正则表示}中的像
为由 $\sum_g\varepsilon(g)e_g$ 张成的直线：符号
表示以重数 $1$ 出现，与
一般理论的要求一致）。
\end{solution}

\begin{solution}{exo:b3:representations:12}
(a) 将\omterm{ex:b3:groups:actions}{共轭类}排列为 $e$ [1]、对换 [6]、
$3$-轮换 [8]、$4$-轮换 [6]、双对换 [3]。
第二个线性\omterm{def:b3:representations:character}{特征标}为 $\chi_2 = \varepsilon$，
其值依次为 $1, -1, 1, -1, 1$。对次数为 $2$ 的
\omterm{def:b3:representations:character}{特征标} $\chi_3$，各列与恒等元所在列的正交性
（$\sum_in_i\chi_i(g) = 0$，当 $g \neq e$），
给出对换类上：$1 - 1 + 2\chi_3 + 3(\chi_4 +
\chi_5) = 0$；借助符号技巧，$\chi_5 = \varepsilon\chi_4$
（次数为 $3$ 的\omterm{def:b3:representations:character}{特征标}乘以线性\omterm{def:b3:representations:character}{特征标}仍为
\omterm{def:b3:representations:rep}{不可约}，范数相同），使 $\chi_4 + \chi_5$ 在
奇类上为零，于是此处 $\chi_3 = 0$。在 $3$-轮换类上：$1 + 1 +
2\chi_3(c_3) + 3(\chi_4 + \chi_5)(c_3) = 0$，且
在偶类上 $\chi_5 = \chi_4$；$c_3$ 列与
自身给出 $\abs{\chi_3}^2$ 的信息，解此小方程组
（并利用 $\chi_3$ 与 $\mathbf 1$ 及
$\varepsilon$ 的行正交性）：$\chi_3 = (2, 0, -1, 0, 2)$，于是 $\chi_4 =
(3, 1, 0, -1, -1)$，及 $\chi_5 = \varepsilon\chi_4 = (3, -1,
0, 1, -1)$。完整的\omterm{def:b3:representations:table}{特征标表}如下：
\begin{center}
\begin{tabular}{c|ccccc}
 & $e$ & $6\,t$ & $8\,c_3$ & $6\,c_4$ & $3\,v$\\
\hline
$\chi_1$ & $1$ & $1$ & $1$ & $1$ & $1$\\
$\chi_2$ & $1$ & $-1$ & $1$ & $-1$ & $1$\\
$\chi_3$ & $2$ & $0$ & $-1$ & $0$ & $2$\\
$\chi_4$ & $3$ & $1$ & $0$ & $-1$ & $-1$\\
$\chi_5$ & $3$ & $-1$ & $0$ & $1$ & $-1$\\
\end{tabular}
\end{center}
（各行范数均为 $1$；各列两两正交：
验证均通过）。

(b) \omterm{ex:b3:groups:actions}{共轭类}数据 $1, 6, 8, 6, 3$ 连同这些次数，恰是
$S_4$ 的数据（轮换型为 $e$、$2$、$3$、$4$、$2{+}2$）。
核：$\ker\chi_2 = \{g : \varepsilon(g) = 1\} = A_4$
（\omterm{ex:b3:groups:actions}{共轭类} $e, c_3, v$：$1 + 8 + 3 = 12$，指数 $2$）；
$\ker\chi_3 = \{g : \chi_3(g) = 2\}$ 为\omterm{ex:b3:groups:actions}{共轭类} $e, v$ 之并，
即克莱因四元群 $V_4$，阶为 $4$，且正规。$\chi_4, \chi_5$
均忠实（$\chi_i(g) = n_i$ 仅在 $e$ 处成立）。核的交
为 $\{e\}$、$V_4$、$A_4$、$G$——
由\cref{exo:b3:representations:6}(b) 可知，这些\emph{全部}
是 $S_4$ 的\omterm{def:b3:groups:normal}{正规子群}。

(c) 满足 $V_4 \subseteq \ker$ 的\omterm{def:b3:representations:character}{特征标}为 $\chi_1,
\chi_2, \chi_3$：它们都经由\omterm{thm:b3:groups:quotient}{商群} $G/V_4$ 分解，它是一个
$6$ 阶群，其\omterm{def:b3:representations:rep}{不可约}次数为 $1, 1, 2$ ——
即 $S_3$ 的\omterm{def:b3:representations:table}{特征标表}。由于\omterm{thm:b3:groups:quotient}{商群}的\omterm{def:b3:representations:table}{特征标表}\emph{是}
两个 $6$ 阶群之间的一个完全不变量（$\Z/6\Z$
将有六个线性\omterm{def:b3:representations:character}{特征标}），故 $G/V_4 \cong
S_3$：\omterm{thm:b3:groups:quotient}{商群}在表中可见，表现为
由核含 $V_4$ 的各行构成的块。
\end{solution}

\begin{solution}{pb:b3:representations:1}
\textbf{1.} 若 $\alpha^n + c_{n-1}\alpha^{n-1} + \dots + c_0 =
0$（$c_i \in \Z$），则 $\alpha^n \in \Z\text{-张成}(1, \dots,
\alpha^{n-1})$，由归纳，每个幂均为这些元的整系数组合，故 $\Z[\alpha]$
由 $1, \alpha, \dots, \alpha^{n-1}$ 生成。反之设
$\Z[\alpha] = \Z g_1 + \dots + \Z g_m$。记 $\alpha g_i =
\sum_j m_{ij}g_j$，其中 $M = (m_{ij}) \in M_m(\Z)$：向量
$g = (g_i)$ 满足 $(\alpha I - M)g = 0$；乘以伴随
矩阵，得 $\det(\alpha I - M)\,g_i = 0$，对每个 $i$，
又因 $1 \in \Z[\alpha]$ 是 $\Z$-组合，
由诸 $g_i$ 而得，故 $\det(\alpha I - M) = 0$：$\alpha$ 是
首一多项式 $\det(XI - M) \in \Z[X]$ 的根。

\textbf{2.} 若 $\Z[\alpha]$ 为 $\alpha$ 的幂（次数至 $n-1$）所张成，
$\Z[\beta]$ 由 $\beta$ 的幂（至 $m - 1$）张成，则
$\Z[\alpha,\beta]$ 由 $nm$ 个乘积
$\alpha^i\beta^j$ 张成（任一单项式均可约化）。子环
$\Z[\alpha + \beta]$ 与 $\Z[\alpha\beta]$ 均为 $\Z$-子模，
即\omterm{def:b3:modules:free}{有限生成} $\Z$-模 $\Z[\alpha, \beta]$ 的子模，
从而\omterm{def:b3:modules:free}{有限生成}（$\Z[\alpha,\beta]$，由
$nm$ 个元生成，是 $\Z^{nm}$ 的像；子模的原像
为 $\Z^{nm}$ 的子模，自由且秩 $\leq nm$（由
\cref{thm:b3:modules:submodule}），其像生成）。由
问题 1，$\alpha + \beta$ 与 $\alpha\beta$ 是代数
整数。

\textbf{3.} 若 $\frac pq$（既约分数）是某个首一
$n$ 次整系数多项式的根，由有理根定理
（去分母：$p^n = -q(\cdots)$）得 $q \mid p^n$，
故 $q = \pm1$。

\textbf{4.} $\chi(g)$ 是若干单位根之和
（见\cref{prop:b3:representations:charbasics}），每项都是代数
整数（$X^N - 1$ 的根）；由问题 2 可得结论。

\textbf{5.} 对 $h \in G$：$\rho(h)S_C\rho(h)^{-1} = \sum_{g\in
C}\rho(hgh^{-1}) = S_C$（$C$ 是一个\omterm{ex:b3:groups:actions}{共轭类}）。由舒尔引理，$S_C =
\omega_C\,\mathrm{id}$；取迹，得 $\abs C\,\chi(g_C) =
\omega_C\, n$。

\textbf{6.} $S_CS_{C'} = \sum_{x \in C, y \in C'}\rho(xy) =
\sum_{z \in G} a(z)\rho(z)$，其中 $a(z) = \#\{(x,y) \in C\times
C' : xy = z\}$。以 $h$ 作共轭，给出 $z$ 的解与
$hzh^{-1}$ 的解之间的双射：$a$ 是\omterm{def:b3:representations:character}{类函数}，取值
于 $\N$，故 $S_CS_{C'} = \sum_{C''}a_{CC'C''}S_{C''}$。
将 $S_C = \omega_C\,\mathrm{id}$ 代入全式，并
比较标量，得
$\omega_C\omega_{C'} = \sum_{C''}a_{CC'C''}\,\omega_{C''}$。

\textbf{7.} 设 $M$ 为 $\Z$-模，由 $1$ 及一切
乘积 $\omega_{C_1}\cdots\omega_{C_k}$ 张成；由问题 6，每个
此类乘积都可约化为 $\Z$-组合，其项为 $1$ 与
$\omega_C$：$M$ \omterm{def:b3:modules:free}{有限生成}，且 $\omega_C M
\subseteq M$（对每个 $C$）。特别地，$\Z[\omega_C]
\subseteq M$ 是\omterm{def:b3:modules:free}{有限生成}的（作为子模，同问题 2），
由问题 1，$\omega_C$ 是代数整数。

\textbf{8.} 对\omterm{def:b3:representations:rep}{不可约}的 $\chi$，其次数为 $n$：
\[
\frac{\abs G}{n} = \frac{\abs G}{n}\langle\chi,\chi\rangle
= \frac1n \sum_{g}\chi(g)\overline{\chi(g)}
= \sum_{C}\frac{\abs C\,\chi(g_C)}{n}\,\overline{\chi(g_C)}
= \sum_C \omega_C\,\overline{\chi(g_C)} .
\]
每个 $\overline{\chi(g_C)} = \chi(g_C^{-1})$ 都是代数
整数（由问题 4），故右端亦为代数整数（问题 2、
7）；它是有理数，从而为整数（由问题 3）：$n \mid \abs
G$。

\textbf{9.} 贝祖：$u\abs C + vn = 1$，其中 $u, v \in \Z$。
于是
\[
\frac{\chi(g_C)}{n} = u\,\frac{\abs C\,\chi(g_C)}{n} +
v\,\chi(g_C) = u\,\omega_C + v\,\chi(g_C),
\]
为代数整数。

\textbf{10.} 设 $0 < \abs{\chi(g_C)} < n$，并令 $\alpha
= \chi(g_C)/n$。所有取值 $\chi(g)$ 都落在 $\Q(\zeta_N)$ 中，$N =
\abs G$（$N$ 次单位根之和）。对 $\sigma \in
\operatorname{Gal}(\Q(\zeta_N)/\Q)$：$\sigma$ 把单位根
映为单位根（$\sigma(\zeta^k) = \zeta^{ak}$，
\cref{thm:b3:galois:cyclotomicirred}），故 $\sigma(\chi(g_C))$
仍是 $n$ 个单位根之和：$\abs{\sigma(\alpha)}
\leq 1$；又 $\sigma(\alpha)$ 是代数整数（它与
$\alpha$ 具有相同的\omterm{def:b3:galois:algebraic}{极小多项式}）。乘积 $P =
\prod_\sigma \sigma(\alpha)$ 在整个\omterm{def:b3:galois:galois}{伽罗瓦群}作用下不变，
从而为有理数（\cref{thm:b3:galois:fundamental}(1)），
且它是代数整数，并满足
\[
0 < \abs P \leq \abs\alpha < 1
\]
（各因子皆非零：$\sigma(\alpha) = 0$ 将迫使 $\alpha =
0$）。这与问题 3 矛盾。故 $\chi(g_C) = 0$ 或
$\abs{\chi(g_C)} = n$，后一情形下 $\rho(g_C)$ 是
数量矩阵（见\cref{prop:b3:representations:charbasics}）。

\textbf{11.} 列正交性（$C \neq \{e\}$）：
$\sum_i \chi_i(e)\overline{\chi_i(g_C)} = 0$，即 $1 +
\sum_{i \geq 2} n_i\overline{\chi_i(g_C)} = 0$。若每个
非平凡的 $\chi_i$（满足 $p \nmid n_i$）在 $g_C$ 处皆为零，则
将其余各项按其因子 $p$ 归并：
\[
-\frac1p = \sum_{i \geq 2,\ p \mid n_i}
\frac{n_i}{p}\,\overline{\chi_i(g_C)},
\]
右端为代数整数——与问题 3 矛盾。故存在某个
非平凡的 $\chi_i$ 满足 $p \nmid n_i$ 且 $\chi_i(g_C) \neq 0$；
由 $\abs C = p^k$ 得 $\gcd(\abs C, n_i) = 1$，再由问题 10
知 $\rho_i(g_C)$ 为数量矩阵。今设 $G$ 为单非阿贝尔群：
$\chi_i$ 为忠实的（见\cref{exo:b3:representations:6}(c)），且
$Z_i = \{g : \rho_i(g) \text{ 为数量矩阵}\}$ 是\omterm{def:b3:groups:normal}{正规子群}
（即数量矩阵在 $\rho_i$ 下的原像；数量矩阵构成
像群中的一个\omterm{def:b3:groups:normal}{正规子群}——事实上是中心子群）。它包含
$g_C \neq e$，故 $Z_i = G$。于是 $\rho_i(G)$ 为阿贝尔群，
且忠实，从而 $G$ 为阿贝尔群：矛盾。\emph{单非阿贝尔群
不含大小为素数幂的\omterm{ex:b3:groups:actions}{共轭类}，
其大小 $> 1$。}

\textbf{12.} 设 $G$ 为阶 $p^aq^b$ 的\omterm{def:b3:groups:simple}{单群}。若 $b = 0$
（即 $G$ 为 $p$-群）：$Z(G) \neq \{e\}$
（见\cref{thm:b3:groups:pfixed}）为\omterm{def:b3:groups:normal}{正规子群}，故 $Z(G) = G$，为阿贝尔群。
否则取 $Q$ 为西罗 $q$-子群，并取 $z \in Z(Q)
\setminus\{e\}$（仍见\cref{thm:b3:groups:pfixed}）。于是 $Q
\subseteq Z_G(z)$，故 $z$ 所在\omterm{ex:b3:groups:actions}{共轭类}的大小为 $[G : Z_G(z)]$
整除 $[G : Q] = p^a$，故为 $p$ 的幂。若大小为 $1$，
则 $z \in Z(G)$：中心是非平凡的\omterm{def:b3:groups:normal}{正规子群}，故
$Z(G) = G$，为阿贝尔群。若大小为 $p^k > 1$：由问题 11
对单非阿贝尔群 $G$ 排除此情形。总之，阶为
$p^aq^b$ 的\omterm{def:b3:groups:simple}{单群}必为阿贝尔群（$\cong \Z/p\Z$）。

\textbf{13.} 对 $\abs G$ 作归纳（$\abs G = 1$ 时可解）。若
$G$ 为\omterm{def:b3:groups:simple}{单群}，由问题 12 知其为阿贝尔群，从而可解。
否则取 $N$ 为非平凡的真\omterm{def:b3:groups:normal}{正规子群}：$\abs
N$ 与 $\abs{G/N}$ 仍具 $p^{a'}q^{b'}$ 的形式，且
阶更小，故由归纳假设 $N$ 与 $G/N$ 皆可解，于是 $G$
可解（见\cref{prop:b3:groups:derived}）。——若有三个素数，
关键一步失效：\omterm{def:b3:groups:sylow}{西罗子群}的指数不再是素数幂，
因而\omterm{def:b3:groups:sylow}{西罗子群}中心元素所在\omterm{ex:b3:groups:actions}{共轭类}的
大小未必为素数幂。某种失效不可避免：
$A_5$，阶为 $2^2\cdot3\cdot5$，是\omterm{def:b3:groups:simple}{单群}且
不可解。

\textbf{14.} \omterm{ex:b3:groups:actions}{共轭类}及其大小见
\cref{exo:b3:groups:11}(a)。$S_5$-共轭类（$c$ 所在者）的大小为
$24$；若它在 $A_5$ 中仍为一个\omterm{ex:b3:groups:actions}{共轭类}，则轨道--稳定子计数
给出 $\abs{Z_{A_5}(c)} = 60/24$，不是整数——
具体地，$Z_{S_5}(c) = \langle c\rangle$ 的阶为 $5$，且
位于 $A_5$ 之中，故 $A_5$-共轭类（$c$ 所在者）的大小为 $60/5 =
12$：$S_5$-共轭类分裂为两个（$c$ 与 $c^2$
仅能借助奇置换在 $S_5$ 中共轭）。

\textbf{15.} 一维\omterm{def:b3:representations:character}{特征标}：次数为 $1$ 的表示
经 $G/D(G)$ 分解，且 $D(A_5) = A_5$（$A_5$ 为单
非阿贝尔群，故 $D(A_5)$ 为正规的非平凡子群——
$A_5$ 非阿贝尔）。于是 $n_1 = 1$，且 $n_2^2 + n_3^2 + n_4^2
+ n_5^2 = 59$，每个 $n_i \geq 2$。可用的平方数有：$4, 9,
16, 25, 36, 49$。四数之和为 $59$ 时，其中最大者
必为 $25$（$36 + 4 + 4 + 9 = 53 < 59$，无法
调整：$36 + 16 + 4 + 4 = 60$，$36 + 9 + 9 + 4 = 58$，$36 +
16 + 9 + 4 = 65$——含 $36$ 或 $49$ 的组合均不可行），
而 $59 - 25 = 34 = 16 + 9 + 9$（唯一可行的分法；$16 + 16 + 4 =
36$，$25 + 9 + 4 = 38$，$25 + 4 + 4 = 33$ 均不成立）：次数为 $1, 3, 3,
4, 5$。

\textbf{16.} $\langle\pi, \mathbf 1\rangle = \frac1{60}(5 +
15\cdot1 + 20\cdot2 + 0 + 0) = 1$（一个\omterm{def:b3:groups:action}{轨道}——伯恩赛德
计数），且 $\langle\pi, \pi\rangle = \frac1{60}(25 + 15 + 80 +
0 + 0) = 2$：$\pi$ 含平凡\omterm{def:b3:representations:character}{特征标}恰一次，其
另一组成部分为一个\omterm{def:b3:representations:rep}{不可约}\omterm{def:b3:representations:character}{特征标}。因此 $\chi_4 = \pi
- \mathbf 1$ \omterm{def:b3:representations:rep}{不可约}，次数为 $4$，取值为 $4, 0,
1, -1, -1$。

\textbf{17.} 对二元子集而言，元素固定 $\{i,j\}$ 当且仅当它
固定或互换 $i, j$：计数依次为 $10$（$e$），$2$（$t =
(1\,2)(3\,4)$ 固定 $\{1,2\}, \{3,4\}$），$1$（$(1\,2\,3)$
固定 $\{4,5\}$），$0, 0$。于是 $\langle\pi_{10},
\pi_{10}\rangle = \frac1{60}(100 + 15\cdot4 + 20\cdot1) = 3$：
三个\omterm{def:b3:representations:rep}{不可约}组成成分，各出现一次。又
$\langle\pi_{10}, \mathbf 1\rangle = \frac1{60}(10 + 30 + 20)
= 1$, $\langle\pi_{10}, \chi_4\rangle = \frac1{60}(40 + 0 +
20\cdot1\cdot1 + 0 + 0) = 1$：故 $\pi_{10} = \mathbf 1 + \chi_4
+ \chi$，其中 $\chi$ \omterm{def:b3:representations:rep}{不可约}，次数为 $10 - 1 - 4 = 5$，且
取值为 $\chi_5 = \pi_{10} - \mathbf 1 - \chi_4 = (5, 1, -1, 0,
0)$。

\textbf{18.} 记 $a, a'$ 为 $\chi_2, \chi_3$ 在
$t$ 处的值，$b, b'$ 为其在 $s = (1\,2\,3)$ 处的值；四者皆为实数（$t$
与 $s$ 均与其逆共轭）。列 $t$ 与
列 $e$ 正交：$1 + 3a + 3a' + 4\cdot0 + 5\cdot1 = 0$，故 $a + a'
= -2$；列 $t$ 与自身作内积：$1 + a^2 + a'^2 + 0 + 1 =
\frac{60}{15} = 4$，故 $a^2 + a'^2 = 2$；于是 $(a + a')^2 = 4
= 2 + 2aa'$，得 $aa' = 1$，且 $a = a' = -1$。列 $s$
与 $e$ 正交：$1 + 3(b + b') + 4 - 5 = 0$，得 $b + b' = 0$；
列 $s$ 与自身作内积：$1 + b^2 + b'^2 + 1 + 1 = \frac{60}
{20} = 3$，得 $b = b' = 0$。

\textbf{19.} 列 $c$ 与列 $e$ 正交：$1 + 3(x + y) +
4(-1) + 5\cdot0 = 0$，故 $x + y = 1$。列 $c$ 与
列 $c^2$（属不同\omterm{ex:b3:groups:actions}{共轭类}，故正交）：$1 + xy + yx + 1
+ 0 = 0$，故 $xy = -1$。于是 $x, y$ 为方程 $T^2 - T - 1 = 0$ 的根：
$\{x, y\} = \{\varphi, \bar\varphi\}$，其中 $\varphi = \frac{1 +
\sqrt5}2$。相容性检验：$x^2 + y^2 = (x+y)^2 - 2xy = 3 =
\frac{60}{12} - 2$——与列自身的恒等式 $1 +
x^2 + y^2 + 1 + 0 = 5$，相符。\omterm{def:b3:representations:table}{特征标表}如下：
\begin{center}
\begin{tabular}{c|ccccc}
 & $e$ & $15\,t$ & $20\,s$ & $12\,c$ & $12\,c^2$\\
\hline
$\chi_1$ & $1$ & $1$ & $1$ & $1$ & $1$\\
$\chi_2$ & $3$ & $-1$ & $0$ & $\varphi$ & $\bar\varphi$\\
$\chi_3$ & $3$ & $-1$ & $0$ & $\bar\varphi$ & $\varphi$\\
$\chi_4$ & $4$ & $0$ & $1$ & $-1$ & $-1$\\
$\chi_5$ & $5$ & $1$ & $-1$ & $0$ & $0$\\
\end{tabular}
\end{center}

\textbf{20.} $\norm{\chi_2}^2 = \frac1{60}\bigl(9 + 15\cdot1 +
0 + 12\varphi^2 + 12\bar\varphi^2\bigr) = \frac{9 + 15 +
36}{60} = 1$，利用 $\varphi^2 + \bar\varphi^2 = (\varphi +
\bar\varphi)^2 - 2\varphi\bar\varphi = 1 + 2 = 3$。同理
$\langle\chi_2, \chi_3\rangle = \frac1{60}(9 + 15 + 0 +
12(2\varphi\bar\varphi)) = \frac{9 + 15 - 24}{60} = 0$。
弗罗贝尼乌斯整除性：$1, 3, 3, 4, 5$ 均整除 $60$。
无理值 $\varphi, \bar\varphi$ 是与 $S_5$ 之间可见的
差异。在 $S_5$ 中，每个元素都与其循环群的
全部生成元共轭（轮换型相同）
——特别地 $c \sim c^2$——迫使\omterm{def:b3:representations:character}{特征标}值为有理数（实际上是
整数）；而在 $A_5$ 中，五阶轮换的分裂
为通向 $\Q(\sqrt5)$ 的道路敞开了大门。

\textbf{21.} \omterm{def:b3:groups:normal}{正规子群} $N$ 是若干\omterm{ex:b3:groups:actions}{共轭类}的并，
包含 $e$，且 $\abs N \mid 60$。候选和如下：
$1{+}15 = 16$, $1{+}20 = 21$, $1{+}12 = 13$, $1{+}24 = 25$,
$1{+}15{+}20 = 36$, $1{+}15{+}12 = 28$, $1{+}15{+}24 = 40$,
$1{+}20{+}12 = 33$, $1{+}20{+}24 = 45$, $1{+}12{+}12 = 25$,
$1{+}15{+}20{+}12 = 48$, $1{+}15{+}20{+}24 = 60 - 12 = \dots$
列出所有含 $1$ 的真部分和：以下各数 $13, 16,
21, 25, 28, 33, 36, 40, 45, 48, 13{+}\dots$ 均不整除 $60$，
除 $1$ 本身外：$N = \{e\}$ 或 $A_5$。单性
由这五个数直接读出。第 11 题的判别准则也清晰可见：
各类大小 $15 = 3\cdot5$、$20 = 4\cdot5$、$12 = 4\cdot3$ 均
为两个素数的乘积——不存在大小为素数幂的类，恰如
伯恩赛德判别准则对\omterm{def:b3:groups:simple}{单群}的要求。

\textbf{22.} $\R^3$ 中转角为 $\theta$ 的旋转，其
特征值为 $1, \eu^{\iu\theta}, \eu^{-\iu\theta}$，迹为 $1 +
2\cos\theta$。当 $\theta = \frac{2\pi}5$ 时，$2\cos\frac{2\pi}5
= \frac{\sqrt5 - 1}2$，从而迹为 $1 + \frac{\sqrt5-1}2 =
\frac{1+\sqrt5}2 = \varphi$。记 $\tau$ 为
$\operatorname{Gal}(\Q(\sqrt5)/\Q)$ 的非平凡元，将其逐项作用于
\omterm{def:b3:representations:table}{特征标表}，可将\omterm{def:b3:representations:character}{特征标}映为\omterm{def:b3:representations:character}{特征标}（它与定义代数可交换：
$\tau\circ\chi$ 为表示的\omterm{def:b3:representations:character}{特征标}，该表示由矩阵
经其元素上的映射得到，
$\tau$ 作用于其元素，或抽象地说，正交关系
均为 $\Q$ 有理的，故 $\tau$ 置换其解）；
$\tau$ 固定 $\chi_1, \chi_4, \chi_5$（取有理值），并
因而必交换 $\chi_2$ 与 $\chi_3$：第二个
$3$ 次\omterm{def:b3:representations:character}{特征标}取共轭值，无需
计算矩阵。几何上，两个表示即
正二十面体作用，及其与 $A_5$ 的一个外
自同构（由一个对换共轭给出）的复合，
这一自同构交换了五阶轮换的两个\omterm{ex:b3:groups:actions}{共轭类}。

\textbf{23.} 对 $z \in Z(G)$，$\rho(z)$ 与每个
$\rho(g)$ 可交换，由舒尔引理得 $\rho(z) = \lambda\,
\mathrm{id}$；由于 $z$ 的阶有限，$\lambda$ 是一个
单位根，且 $\abs{\chi(z)} = \abs\lambda\,n = n$。于是
\[
\abs G = \abs G\,\langle\chi, \chi\rangle
= \sum_{g \in G}\abs{\chi(g)}^2
\geq \sum_{z \in Z(G)}\abs{\chi(z)}^2
= \abs{Z(G)}\,n^2,
\]
即 $n^2 \leq [G : Z(G)]$。阶为 $8$ 的
非阿贝尔群（$D_4$ 或 $Q_8$）有五个\omterm{ex:b3:groups:actions}{共轭类}，故有五个
\omterm{def:b3:representations:rep}{不可约}\omterm{def:b3:representations:character}{特征标}的次数，满足 $\sum n_i^2 = 8$；将 $8$
写成五个 $\geq 1$ 的平方和只有一种方式：$1 + 1 + 1 + 1
+ 4$：次数为 $1, 1, 1, 1, 2$。两群的中心阶均为
阶为 $2$，且 $2$ 次\omterm{def:b3:representations:character}{特征标}取等号：
$2^2 = 4 = [G : Z(G)]$——此界是紧的。对 $A_5$，
$Z = \{e\}$，界式为 $n^2 \leq 60$：$1, 3, 3, 4, 5$ 均满足，
且尚有余量（$25 \leq 60$），这是必然的，
因为若取等号，由同一推导链将迫使 $\chi$
在中心之外为零。

\textbf{24.} $\rho(g)^m = \mathrm{id}$，其中 $m$ 为
$g$ 的阶，故 $\rho(g)$ 被 $X^m - 1$ 零化，该多项式在
$\C$ 上分裂且根皆为单根：可对角化，特征值为
$\lambda_1, \dots, \lambda_n$ 为单位根，在某个
特征基 $(e_i)$ 中。乘积 $e_ie_j$（$i \leq j$）构成
$\operatorname{Sym}^2V$ 的特征基，其特征值为
$\lambda_i\lambda_j$，而 $e_i \wedge e_j$（$i < j$）构成
$\Lambda^2V$ 的特征基；由于
\[
\sum_{i<j}\lambda_i\lambda_j
= \frac{(\sum_i\lambda_i)^2 - \sum_i\lambda_i^2}2
= \frac{\chi(g)^2 - \chi(g^2)}2,
\]
而对称和则是加上 $\sum_i\lambda_i^2$ 而非
减去它，两式均由此得出。对 $\chi_2 = (3, -1,
0, \varphi, \bar\varphi)$ 在类 $(e, (2,2)\text{-},
3\text{-轮换}, c, c^2)$：平方将双
对换映为 $e$，三阶轮换映为三阶轮换，
$c$ 所在的类映到 $c^2$ 所在的类，反之亦然（$c^{-1} \sim
c$ 在 $A_5$ 中成立，故 $c^4 \sim c$）。于是 $\chi_2(g^2)$ 为 $(3,
3, 0, \bar\varphi, \varphi)$，并利用 $\varphi^2 = \varphi
+ 1$，$\bar\varphi = 1 - \varphi$：
\[
\Lambda^2\chi_2 = (3, -1, 0, \varphi, \bar\varphi) = \chi_2,
\qquad
\operatorname{Sym}^2\chi_2 = (6, 2, 0, 1, 1) .
\]
将后者分解，类大小依次为 $1, 15, 20, 12,
12$：$\langle\cdot, \mathbf 1\rangle = \frac1{60}(6 +
15\cdot2 + 0 + 12 + 12) = 1$；$\langle\cdot, \chi_5\rangle =
\frac1{60}(30 + 30) = 1$；$\langle\cdot, \chi_4\rangle =
\frac1{60}(24 - 12 - 12) = 0$；$\langle\cdot, \chi_2\rangle =
\frac1{60}(18 - 30 + 12(\varphi + \bar\varphi)) = 0$，且
对 $\chi_3$ 同理。故 $\operatorname{Sym}^2\chi_2 =
\mathbf 1 + \chi_5$（维数 $6 = 1 + 5$），而
$\chi_2\otimes\chi_2 = \mathbf 1 + \chi_2 + \chi_5$
（维数 $9 = 1 + 3 + 5$）。几何上，等变
同构 $\Lambda^2\R^3 \to \R^3$，$u \wedge v \mapsto u
\times v$，对旋转群而言恰为 $\Lambda^2\chi_2 = \chi_2$；
对称平方中的直和因子 $\mathbf 1$ 是
不变二次型 $x^2 + y^2 + z^2$，而
$\chi_5$ 作用于五维的无迹对称张量
空间上（即调和二次式）。

\textbf{25.} $c$ 列与 $c^2$ 列的内积：
\[
1\cdot1 + \varphi\bar\varphi + \bar\varphi\varphi +
(-1)(-1) + 0 = 1 - 1 - 1 + 1 + 0 = 0,
\]
正如正交性对不同类所要求的（$\varphi
\bar\varphi = -1$）。$c$ 列与其自身作内积：$1 +
\varphi^2 + \bar\varphi^2 + 1 + 0 = 1 + 3 + 1 = 5 = 60/12 =
\abs{Z_{A_5}(c)}$。正则\omterm{def:b3:representations:character}{特征标} $\sum_in_i\chi_i$ 在
四个非单位元所在列上：
\[
1 - 3 - 3 + 0 + 5 = 0, \qquad
1 + 0 + 0 + 4 - 5 = 0,
\]
在双对换与三阶轮换所在类上，以及在
$c$ 所在的类上（$c^2$ 所在的类是其伽罗瓦共轭）：
\[
1 + 3\varphi + 3\bar\varphi - 4 + 0 = 1 + 3 - 4 = 0,
\]
利用 $\varphi + \bar\varphi = 1$。该表通过了
每项检验：它正是 $A_5$ 的\omterm{def:b3:representations:table}{特征标表}。
\end{solution}
